
\documentclass{article}
\usepackage{colortbl}
\usepackage{makecell}
\usepackage{multirow}
\usepackage{supertabular}

\begin{document}

\newcounter{utterance}

\centering \large Interaction Transcript for game `cladder', experiment `full\_v1.5\_default', episode 8398 with Dialogue Architects\_\_qwen3.5{-}9b\_\_SCoRe\_Qwen3.5{-}9B.
\vspace{24pt}

{ \footnotesize  \setcounter{utterance}{1}
\setlength{\tabcolsep}{0pt}
\begin{supertabular}{c@{$\;$}|p{.15\linewidth}@{}p{.15\linewidth}p{.15\linewidth}p{.15\linewidth}p{.15\linewidth}p{.15\linewidth}}
    \# & \multicolumn{2}{c}{Player} && \multicolumn{2}{c}{Game Master} \\
    \hline

    \theutterance \stepcounter{utterance}  
    & & & \multicolumn{4}{p{0.6\linewidth}}{
        \cellcolor[rgb]{0.9,0.9,0.9}{
            \makecell[{{p{\linewidth}}}]{
                \texttt{\tiny{[P1$\langle$GM]}}
                \texttt{You are an expert in causal inference. The following question is not a typical commonsense query, but rather a meticulously designed question created by a professor specializing in causal inference, intended to assess the students' mastery of the course content.} \\
\texttt{Imagine a self{-}contained, hypothetical world with only the following conditions, and without any unmentioned factors or causal relationships: Jyka has a direct effect on hwax. Hwax has a direct effect on lirg. The overall probability of jyka is 14\%. The probability of not jyka and lirg is 72\%. The probability of jyka and lirg is 12\%. Is the chance of lirg smaller when observing jyka?} \\
\texttt{Start your answer with 'Yes' or 'No', followed by additional reasoning or evidence to support your explanation.} \\
            }
        }
    }
     \\ \\

    \theutterance \stepcounter{utterance}  
    & \multicolumn{4}{p{0.6\linewidth}}{
        \cellcolor[rgb]{0.561,0.737,0.561}{
            \makecell[{{p{\linewidth}}}]{
                \texttt{\tiny{[P1$\rangle$GM]}}
                \texttt{No} \\
\\ 
\texttt{**Reasoning:**} \\
\\ 
\texttt{To determine if the chance of *lirg* is smaller when observing *jyka*, we must compare the marginal probability of *lirg* (\$P(\textbackslash{}text\{lirg\})\$) with the conditional probability of *lirg* given *jyka* (\$P(\textbackslash{}text\{lirg\} | \textbackslash{}text\{jyka\})\$).} \\
\\ 
\texttt{1.  **Calculate the marginal probability of *lirg* (\$P(\textbackslash{}text\{lirg\})\$):**} \\
\texttt{    We are given the probability of the joint event where *jyka* is false and *lirg* is false: \$P(\textbackslash{}neg \textbackslash{}text\{jyka\} \textbackslash{}cap \textbackslash{}neg \textbackslash{}text\{lirg\}) = 0.72\$.} \\
\texttt{    Using the complement rule, the probability that *lirg* occurs is:} \\
\texttt{    \$\$P(\textbackslash{}text\{lirg\}) = 1 {-} P(\textbackslash{}neg \textbackslash{}text\{lirg\})\$\$} \\
\texttt{    Since the events are binary, \$P(\textbackslash{}neg \textbackslash{}text\{lirg\}) = P(\textbackslash{}neg \textbackslash{}text\{jyka\} \textbackslash{}cap \textbackslash{}neg \textbackslash{}text\{lirg\}) + P(\textbackslash{}text\{jyka\} \textbackslash{}cap \textbackslash{}neg \textbackslash{}text\{lirg\})\$. However, we can more directly use the total probability space.} \\
\texttt{    We know \$P(\textbackslash{}text\{jyka\} \textbackslash{}cap \textbackslash{}text\{lirg\}) = 0.12\$.} \\
\texttt{    We know \$P(\textbackslash{}neg \textbackslash{}text\{jyka\} \textbackslash{}cap \textbackslash{}neg \textbackslash{}text\{lirg\}) = 0.72\$.} \\
\texttt{    The sum of these two disjoint events is \$0.12 + 0.72 = 0.84\$.} \\
\texttt{    The remaining probability mass corresponds to the cases where *lirg* is true but *jyka* is false (\$P(\textbackslash{}neg \textbackslash{}text\{jyka\} \textbackslash{}cap \textbackslash{}text\{lirg\})\$) or where *jyka* is true and *lirg* is false (\$P(\textbackslash{}text\{jyka\} \textbackslash{}cap \textbackslash{}neg \textbackslash{}text\{lirg\})\$).} \\
\\ 
\texttt{    Let's calculate \$P(\textbackslash{}text\{lirg\})\$ directly by summing the probabilities of all scenarios where *lirg* is true:} \\
\texttt{    \$\$P(\textbackslash{}text\{lirg\}) = P(\textbackslash{}text\{jyka\} \textbackslash{}cap \textbackslash{}text\{lirg\}) + P(\textbackslash{}neg \textbackslash{}text\{jyka\} \textbackslash{}cap \textbackslash{}text\{lirg\})\$\$} \\
\texttt{    We know \$P(\textbackslash{}text\{jyka\}) = 0.14\$. Since \$P(\textbackslash{}text\{jyka\}) = P(\textbackslash{}text\{jyka\} \textbackslash{}cap \textbackslash{}text\{lirg\}) + P(\textbackslash{}text\{jyka\} \textbackslash{}cap \textbackslash{}neg \textbackslash{}text\{lirg\})\$, we can find \$P(\textbackslash{}text\{jyka\} \textbackslash{}cap \textbackslash{}neg \textbackslash{}text\{lirg\})\$:} \\
\texttt{    \$\$P(\textbackslash{}text\{jyka\} \textbackslash{}cap \textbackslash{}neg \textbackslash{}text\{lirg\}) = P(\textbackslash{}text\{jyka\}) {-} P(\textbackslash{}text\{jyka\} \textbackslash{}cap \textbackslash{}text\{lirg\}) = 0.14 {-} 0.12 = 0.02\$\$} \\
\\ 
\texttt{    Now, let's find \$P(\textbackslash{}neg \textbackslash{}text\{jyka\} \textbackslash{}cap \textbackslash{}text\{lirg\})\$. We know the total probability of \$\textbackslash{}neg \textbackslash{}text\{jyka\}\$ is \$1 {-} 0.14 = 0.86\$.} \\
\texttt{    \$\$P(\textbackslash{}neg \textbackslash{}text\{jyka\}) = P(\textbackslash{}neg \textbackslash{}text\{jyka\} \textbackslash{}cap \textbackslash{}text\{lirg\}) + P(\textbackslash{}neg \textbackslash{}text\{jyka\} \textbackslash{}cap \textbackslash{}neg \textbackslash{}text\{lirg\})\$\$} \\
\texttt{    \$\$0.86 = P(\textbackslash{}neg \textbackslash{}text\{jyka\} \textbackslash{}cap \textbackslash{}text\{lirg\}) + 0.72\$\$} \\
\texttt{    \$\$P(\textbackslash{}neg \textbackslash{}text\{jyka\} \textbackslash{}cap \textbackslash{}text\{lirg\}) = 0.86 {-} 0.72 = 0.14\$\$} \\
\\ 
\texttt{    Now sum the parts where *lirg* is true:} \\
\texttt{    \$\$P(\textbackslash{}text\{lirg\}) = 0.12 (\textbackslash{}text\{from jyka\}) + 0.14 (\textbackslash{}text\{from not jyka\}) = 0.26\$\$} \\
\\ 
\texttt{2.  **Calculate the conditional probability \$P(\textbackslash{}text\{lirg\} | \textbackslash{}text\{jyka\})\$:**} \\
\texttt{    \$\$P(\textbackslash{}text\{lirg\} | \textbackslash{}text\{jyka\}) = \textbackslash{}frac\{P(\textbackslash{}text\{jyka\} \textbackslash{}cap \textbackslash{}text\{lirg\})\}\{P(\textbackslash{}text\{jyka\})\}\$\$} \\
\texttt{    \$\$P(\textbackslash{}text\{lirg\} | \textbackslash{}text\{jyka\}) = \textbackslash{}frac\{0.12\}\{0.14\} \textbackslash{}approx 0.857\$\$} \\
\\ 
\texttt{3.  **Comparison:**} \\
\texttt{    The question asks: "Is the chance of *lirg* smaller when observing *jyka*?"} \\
\texttt{    This translates to: Is \$P(\textbackslash{}text\{lirg\} | \textbackslash{}text\{jyka\}) < P(\textbackslash{}text\{lirg\})\$?} \\
\texttt{    Comparing the values: \$0.857\$ (conditional) vs \$0.26\$ (marginal).} \\
\texttt{    \$0.857 > 0.26\$.} \\
\\ 
\texttt{The probability of *lirg* increases significantly when *jyka* is observed (from 26\% to \textasciitilde{}86\%). This is consistent with the causal structure described (*jyka* \$\textbackslash{}to\$ *hwax* \$\textbackslash{}to\$ *lirg*), where *jyka* is a positive cause of *lirg*. Therefore, the chance is not smaller; it is larger.} \\
\\ 
\texttt{**Answer:** No} \\
            }
        }
    }
    & & \\ \\

    \theutterance \stepcounter{utterance}  
    & & & \multicolumn{2}{p{0.3\linewidth}}{
        \cellcolor[rgb]{0.9,0.9,0.9}{
            \makecell[{{p{\linewidth}}}]{
                \texttt{\tiny{[GM$|$GM]}}
                \texttt{No} \\
            }
        }
    }
    & & \\ \\

    \theutterance \stepcounter{utterance}  
    & & & \multicolumn{2}{p{0.3\linewidth}}{
        \cellcolor[rgb]{0.9,0.9,0.9}{
            \makecell[{{p{\linewidth}}}]{
                \texttt{\tiny{[GM$|$GM]}}
                \texttt{no} \\
            }
        }
    }
    & & \\ \\

    \theutterance \stepcounter{utterance}  
    & & & \multicolumn{2}{p{0.3\linewidth}}{
        \cellcolor[rgb]{0.9,0.9,0.9}{
            \makecell[{{p{\linewidth}}}]{
                \texttt{\tiny{[GM$|$GM]}}
                \texttt{game\_result = WIN} \\
            }
        }
    }
    & & \\ \\

\end{supertabular}
}

\end{document}
